\section{Expose a complete baseline block emitter}
Implement \texttt{h263\_emit\_coefficients}. Pass an eight-byte-aligned mapped
32-byte request in \texttt{\$a0}. Its unsigned native-endian doublewords are:
\begin{center}
  \tagpdfsetup{table/header-rows={1}}
  \begin{tabular}{rll}
    \textbf{Offset} & \textbf{Field}         & \textbf{Meaning}                             \\
    0               & Level pointer          & Two-byte-aligned, 64 natural-order halfwords \\
    8               & Writer-context pointer & Eight-byte-aligned prior 24-byte context     \\
    16              & Block type             & 0 INTER; 1 INTRA                             \\
    24              & Reserved               & Zero
  \end{tabular}
\end{center}
All ordinary levels lie in $[-127,127]$, with zero omitted from events. For INTRA,
index zero instead contains a positive DC code from one through 255, excluding
128. There is no quantizer parameter here: these are already selected levels.
The emitter must preserve their exact meaning, not reconstruct coefficients or
perform another quantization.

Request, context and level input are disjoint; the byte output arena is disjoint
from all three. All complete extents are mapped and non-wrapping. The writer
context retains its earlier capacity bound of one MiB and a used count within
that capacity. Preserve saved general registers, incoming stack pointer, request,
level bytes and context pointer/capacity fields. Supply an aligned stack with
private space; vector registers are caller-saved.

Validate metadata before level reads. A malformed request, context or level
returns status one and count zero. Reads to find an invalid level are permitted
after valid metadata. No caller output access or state change is permitted on
such rejection. A chosen plan that exceeds remaining capacity returns status two
and count zero, also before caller-buffer access. A failed bit-writer dependency
returns status three and count zero; prior caller-buffer reads needed for private
staging are permitted, but no caller output or count may be written.

On success, append one complete baseline block payload at the existing bit cursor,
preserve every bit outside it, and commit the new used count. Return status zero
and the number of appended bits in \texttt{\$v1}. An empty INTER block appends
nothing and makes no caller-buffer access or context write. An INTRA block always
appends its DC byte, even when every ordinary AC level is zero. This function
does not emit a macroblock header, coded-block pattern, alignment padding or file.

Use \path{construction/h263-coefficient-construction/main.asm}. Its supplied scan and
complete compact-code tables are numeric data. Replace the emitter stub and carry
in your completed \texttt{h263\_write\_bits} entry and helpers, replacing its stub.
The default INTRA observation has DC code 255, level one at natural index one,
level minus one at natural index eight and zero elsewhere, beginning at bit offset
three in a buffer of \texttt{A5} bytes.
\begin{verbatim}
python3 construction/h263-coefficient-construction/prepare.py \
    --source my-coefficients.asm --intra 1 --pattern worked \
    --capacity 192 --used 3 --output worked-coefficients.asm
\end{verbatim}
Preparation changes the provided level data and typed controls, retaining all
instructions and code tables. Existing output files are refused. The completed
compact reference appends sixteen bits and prints sixteen; the valid escape-only
alternative appends fifty-two and prints fifty-two. The prefix is retained in
both cases, and both exit zero. An unfinished emitter reports development status
four and exits one.

At \texttt{coeff\_caller\_return}, before the first caller instruction, output
and context are complete and return registers contain status and appended bits.
Named status/bit-count fields have not yet received their copies. Inspect the
committed bit interval and decode it; a differing printed bit count alone is
not evidence of a wrong coefficient block.
